% Part: proof-theory
% Chapter: sequent-calculus
% Section: proof-examples

\documentclass[../../../include/open-logic-section]{subfiles}

\begin{document}

\olfileid[hi]{pt}{seq}{ex}

\olsection{प्रमाणों के उदाहरण}

अनुक्रमों के !!{proof}s देने के लिए सामान्यतः अंतिम अनुक्रम से ऊपर की ओर काम
करना सुविधाजनक होता है: उसमें किसी !!a{formula} को सक्रिय !!{formula} चुनें,
अनुक्रम के ऊपर अनुरूपी नियम के आधार-वाक्य लिखें, और अभिगृहीतियों तक पहुँचने
तक यही प्रक्रिया दोहराएँ। उदाहरणतः, मान लें कि हमें यह अनुक्रम सिद्ध करना है:
\[(!C \land !D) \lif !E \Sequent !C \lif (!D \lif !E)\]
$!C \lif (!D \lif !E)$ पर विचार करें: यह $!A \lif !B$ रूप का है और अनुक्रम
की दाईं ओर आता है। पूरे अनुक्रम का \Log{G1c} के $\RightR{\lif}$ नियम से
मिलान करने पर,
\[\Axiom$ !A, \Gamma \fCenter \Delta, !B$
\RightLabel{\RightR{\lif}}
\UnaryInf$ \Gamma \fCenter \Delta, !A \lif !B $
\DisplayProof\]
यह संकेत मिलता है कि हमें $\Gamma = \{!C \lif (!D \lif !E)\}$, $\Delta =
\emptyset$, $!A \ident !C$ और $!B \ident !D \lif !E$ लेना चाहिए। अतः प्रमाण
का अंतिम निगमन यह हो सकता है:
\[\Axiom$ !C, (!C \land !D) \lif !E \fCenter !D \lif !E$
\RightLabel{\RightR{\lif}}
\UnaryInf$(!C \land !D) \lif !E \fCenter !C \lif (!D \lif !E)$
\DisplayProof\]
यदि हम $!D \lif !E$ को सक्रिय सूत्र मानकर यही विचार दोहराएँ, तो मिलता है:
\[\Axiom$ !D, !C, (!C \land !D) \lif !E \fCenter !E$
\RightLabel{\RightR{\lif}}
\UnaryInf$!C, (!C \land !D) \lif !E \fCenter !D \lif !E$
\RightLabel{\RightR{\lif}}
\UnaryInf$(!C \land !D) \lif !E \fCenter !C \lif (!D \lif !E)$
\DisplayProof\]
अब हमें $(!C \land !D) \lif !E$ पर ध्यान देकर \Log{G1c} का
$\LeftR{\lif}$ नियम उपयोग करना है,
\[\Axiom$ \Gamma \fCenter \Delta, !A$
\Axiom$ !B, \Gamma \fCenter \Delta$
\RightLabel{\LeftR{\lif}}
\BinaryInf$ !A \lif !B, \Gamma \fCenter \Delta$
\DisplayProof\]
इससे हमें मिलता है:
\[\Axiom$!D, !C \fCenter !E, !C \land !D$
\Axiom$!D, !C, !E \fCenter !E$
\RightLabel{\LeftR{\lif}}
\BinaryInf$ !D, !C, (!C \land !D) \lif !E \fCenter !E$
\RightLabel{\RightR{\lif}}
\UnaryInf$!C, (!C \land !D) \lif !E \fCenter !D \lif !E$
\RightLabel{\RightR{\lif}}
\UnaryInf$(!C \land !D) \lif !E \fCenter !C \lif (!D \lif !E)$
\DisplayProof\]
ऊपर के बाएँ अनुक्रम में !!{formula} $!C \land !D$ को प्रमुख मानकर
$\RightR{\land}$ नियम उपयोग करने पर हमें मिलता है:
\[
\Axiom$!D, !C \fCenter !E, !C$
\Axiom$!D, !C \fCenter !E, !D$
\RightLabel{\RightR{\land}}
\BinaryInf$!D, !C \fCenter !E, !C \land !D$
\Axiom$!D, !C, !E \fCenter !E$
\RightLabel{\RightR{\lif}}
\BinaryInf$ !D, !C, (!C \land !D) \lif !E \fCenter !E$
\RightLabel{\RightR{\lif}}
\UnaryInf$!C, (!C \land !D) \lif !E \fCenter !D \lif !E$
\RightLabel{\RightR{\lif}}
\UnaryInf$(!C \land !D) \lif !E \fCenter !C \lif (!D \lif !E)$
\DisplayProof
\]
यहाँ तक सब ठीक है। ध्यान दें कि हमने जो कुछ किया वह \Log{G3c} में भी काम
करता है, क्योंकि अब तक उपयोग किए गए नियम \Log{G1c} और~\Log{G3c} में समान
हैं। हम !!a{proof} के ऐसे खंड पर पहुँचे हैं जिसमें ऊपर के सभी अनुक्रमों में
बाईं और दाईं ओर समान !!{formula} है। किंतु वे अभी ठीक-ठीक अभिगृहीतियाँ नहीं हैं।

\Log{G1c} में हमें ऊपर के अनुक्रमों के !!{proof}s, $!A \Sequent !A$ रूप की
अभिगृहीतियों से देने होंगे। दुर्बलीकरण नियमों से यह सरलता से होता है; उदाहरणतः
सबसे ऊपर बाएँ अनुक्रम $!D, !C \Sequent !E, !C$ को इस प्रकार !!{prove} किया
जा सकता है:
\[
\Axiom$!C \fCenter !C$
\RightLabel{\LeftR{\Weakening}}
\UnaryInf$!D, !C \fCenter !C$
\RightLabel{\RightR{\Weakening}}
\UnaryInf$!D, !C \fCenter !E, !C$
\]
कुल मिलाकर, \Log{G1c} में हमारा !!a{proof} इस प्रकार दिखता है:
\[
\Axiom$!C \fCenter !C$
\RightLabel{\LeftR{\Weakening}}
\UnaryInf$!D, !C \fCenter !C$
\RightLabel{\RightR{\Weakening}}
\UnaryInf$!D, !C \fCenter !E, !C$
\Axiom$!D \fCenter !D$
\RightLabel{\LeftR{\Weakening}}
\UnaryInf$!D, !C \fCenter !D$
\RightLabel{\RightR{\Weakening}}
\UnaryInf$!D, !C \fCenter !E, !D$
\RightLabel{\RightR{\land}}
\BinaryInf$!D, !C \fCenter !E, !C \land !D$
\Axiom$!E \fCenter !E$
\RightLabel{\LeftR{\Weakening}}
\UnaryInf$!C, !E \fCenter !C$
\RightLabel{\LeftR{\Weakening}}
\UnaryInf$!D, !C, !E \fCenter !E$
\RightLabel{\LeftR{\lif}}
\BinaryInf$ !D, !C, (!C \land !D) \lif !E \fCenter !E$
\RightLabel{\RightR{\lif}}
\UnaryInf$!C, (!C \land !D) \lif !E \fCenter !D \lif !E$
\RightLabel{\RightR{\lif}}
\UnaryInf$(!C \land !D) \lif !E \fCenter !C \lif (!D \lif !E)$
\DisplayProof
\]
इस प्रमाण की संरचना (विशेषतः इसकी ऊँचाई) इस बात पर निर्भर नहीं करती कि
!!{formula}s $!C$, $!D$ और~$!E$ क्या हैं। ऊपर के !!{proof} में इन चिह्नों को
किन्हीं भी !!{formula}s से बदल सकते हैं और परिणाम फिर भी \Log{G1c} का
!!{proof} रहेगा। \Log{G1c} का यह प्रमाण \emph{योजनाबद्ध} है।

\Log{G3c} में अभिगृहीतियाँ $!A, \Gamma \Sequent \Delta, !A$ रूप के अनुक्रम
हैं, जहाँ $!A$ परमाण्विक होना चाहिए। अतः यद्यपि $!D, !C \Sequent !E, !C$,
\Log{G3c} की अभिगृहीति जैसा दिखता है ($\Gamma = \{!D\}$; $\Delta = \{!E\}$),
वह \emph{वास्तव में} तभी अभिगृहीति है जब $!C$ सचमुच परमाण्विक हो। हमारा
प्रमाण-खंड \Log{G3c} में तभी पूरा प्रमाण है जब $!C$, $!D$ और $!E$ परमाण्विक
!!{formula}s हों।

यद्यपि कठोर अर्थ में हमने यह नहीं दिखाया कि
\[\Log{G1c} \Proves (!C \land !D) \lif !E \fCenter !C \lif (!D \lif
!E),\] फिर भी व्यवहार में हमें इतना ही करना है (और हम इतना ही कर सकते हैं)।
कारण यह है कि $!A, \Gamma \Sequent \Delta, !A$ रूप का हर अनुक्रम \Log{G3c}
में सिद्ध किया जा सकता है, भले ही स्वयं $!A$ परमाण्विक न हो। इसे हम $!A$ की
गहराई~$\depth{!A}$ पर आगमन द्वारा दिखा सकते हैं।

\begin{defn}\ollabel{defn:depth}
!!a{formula} की \emph{गहराई}~$\depth{!A}$ इस प्रकार परिभाषित है:
\begin{enumerate}
  \item यदि $!A$ परमाण्विक है, तो $\depth{!A} = 0$।
  \item यदि $!A \ident \lnot !B$, $!A \ident \lforall[x][A]$, अथवा
  $!A \ident \lexists[x][A]$, तो $\depth{!A} = \depth{!B} + 1$।
  \item यदि $!A \ident (!B \land !C)$, $(!B \lor !C)$, अथवा $(!B \lif
  !C)$, तो $\depth{!A} = \max(\depth{!B},\depth{!C}) + 1$।
\end{enumerate}
\end{defn}

यह सिद्ध करना कठिन नहीं कि किसी भी पद~$t$ के लिए $\depth{A(x)} =
\depth{A(t)}$। हमारे लिए महत्त्वपूर्ण तथ्य यह है कि किसी भी नियम में प्रमुख
!!{formula} की गहराई सक्रिय !!{formula}s की गहराई से सदैव अधिक होती है।
उदाहरणतः हमारे पास है:
\begin{align*}
\depth{P(x)} = \depth{P(a)} & = 0,\\
\depth{\lforall[x][P(x)]} & = 1,\\
\depth{(\lforall[x][P(x)] \lif P(a))} & = \max(1,0) + 1 = 2.
\end{align*}
इसका उपयोग हम~$\depth{!A}$ पर आगमन द्वारा कुछ परिणाम सिद्ध करने में कर सकते
हैं, जैसे निम्न परिणाम।

\begin{prop}\ollabel{prop:G1c-axioms}
किसी भी~$!A$ के लिए $!A \Sequent !A$ का $\Log{G1c}$ में ऐसा !!a{proof} है
जिसमें सभी अभिगृहीतियाँ परमाण्विक हैं।
\end{prop}

\begin{proof}
  ~$\depth{!A}$ पर आगमन द्वारा। यदि $\depth{!A} = 0$, तो $!A$ परमाण्विक है
  और $!A \Sequent !A$ पहले से अपेक्षित रूप का \Log{G1c}-!!{proof} है। यदि
  $\depth{!A} > 0$, तो $!A$ कम गहराई वाले !!{formula}s से बना है, जैसे
  $!A \ident \lnot !B$, $!A \ident (!B \land !C)$, अथवा $!A \ident
  \lforall[x][A(x)]$। आगमन परिकल्पना है कि $\depth{!D} < \depth{!A}$ वाले
  सभी $!D$ के लिए परमाण्विक अभिगृहीतियों से $\Log{G1c} \Proves !D \Sequent
  !D$।

  यदि $!A \ident \lnot !B$, तो $\depth{!B} < \depth{!A}$ और आगमन परिकल्पना
  से परमाण्विक अभिगृहीतियों से $!B \Sequent !B$ का \Log{G1c} में कोई
  !!a{proof}~$\pi_1$ है। इसे हम $\lnot !B \Sequent \lnot !B$ के !!a{proof}
  तक इस प्रकार बढ़ा सकते हैं:
  \[
  \AxiomC{}
  \RightLabel{$\pi_1$}
  \Deduce$!B \fCenter !B$
  \RightLabel{\LeftR{\lnot}}
  \UnaryInf$\lnot !B, !B \fCenter $
  \RightLabel{\RightR{\lnot}}
  \UnaryInf$\lnot !B \fCenter \lnot !B$
  \DisplayProof
\]

  यदि $!A \ident (!B \land !C)$, तो $\depth{!B} < \depth{!A}$ और
  $\depth{!C} < \depth{!A}$। आगमन परिकल्पना से परमाण्विक अभिगृहीतियों से
  $!B \Sequent !B$ और $!C \Sequent !C$ के \Log{G1c} में !!{proof}s~$\pi_1$
  और~$\pi_2$ हैं। इन्हें हम $!B \land !C \Sequent !B \land !C$ के !!a{proof}
  तक इस प्रकार बढ़ा सकते हैं:
  \[
  \AxiomC{}
  \RightLabel{$\pi_1$}
  \Deduce$!B, \fCenter !B$
  \RightLabel{\LeftR{\Weakening}}
  \UnaryInf$!B, !C, \fCenter !B$
  \RightLabel{\LeftR{\land}}
  \UnaryInf$!B \land !C, !C\fCenter !B$
  \RightLabel{\LeftR{\land}}
  \UnaryInf$!B \land !C, !B \land !C \fCenter !B$
  \AxiomC{}
  \RightLabel{$\pi_2$}
  \Deduce$!C, \fCenter !C$
  \RightLabel{\LeftR{\Weakening}}
  \UnaryInf$!B, !C, \fCenter !C$
  \RightLabel{\LeftR{\land}}
  \UnaryInf$!B \land !C, !C \fCenter !C$
  \RightLabel{\LeftR{\land}}
  \UnaryInf$!B \land !C, !B \land !C \fCenter !C$
  \RightLabel{\RightR{\land}}
  \BinaryInf$!B \land !C \fCenter !B \land !C$
  \RightLabel{\LeftR{\Contraction}}
  \UnaryInf$!B \land !C \fCenter !B$
  \DisplayProof
  \]

  अब मान लें कि $!A \ident \lforall[x][B(x)]$। ऐसा !!a{constant}~$c$ लें जो
  $\lforall[x][B(x)]$ में न आता हो। तब $\depth{B(c)} = \depth{!B(x)} <
  \depth{!A}$, और आगमन परिकल्पना से परमाण्विक अभिगृहीतियों से $!B(c)
  \Sequent !B(c)$ का \Log{G1c} में कोई !!a{proof}~$\pi_1$ है। इसे हम
  $\lforall[x][B(x)] \Sequent \lforall[x][B(x)]$ के !!a{proof} तक इस प्रकार
  बढ़ा सकते हैं:
  \[
  \AxiomC{}
  \RightLabel{$\pi_1$}
  \Deduce$!B(c) \fCenter !B(c)$
  \RightLabel{\LeftR{\lforall}}
  \UnaryInf$\lforall[x][!B(x)] \fCenter !B(c)$
  \RightLabel{\RightR{\lforall}}
  \UnaryInf$\lforall[x][!B(x)] \fCenter \lforall[x][!B(x)]$
  \DisplayProof
\]

  शेष स्थितियाँ अभ्यास के लिए छोड़ी गई हैं।
\end{proof}

\begin{prob}
  $!A \ident (!B \lor !C)$, $!A \ident (!B \lif !C)$ और $!A \ident
  \lexists[x][B(x)]$ की स्थितियाँ पूरा करके \olref{prop:G1c-axioms} का प्रमाण
  पूर्ण करें।
\end{prob}

अपने प्रयोजन के लिए, $!A \ident \lnot !B$ की स्थिति में हम यह !!{proof} भी
उपयोग कर सकते थे:
\[
\AxiomC{}
\RightLabel{$\pi_1$}
\Deduce$!B \fCenter !B$
\RightLabel{\RightR{\lnot}}
\UnaryInf$\fCenter !B, \lnot !B$
\RightLabel{\LeftR{\lnot}}
\UnaryInf$\lnot !B \fCenter \lnot !B$
\DisplayProof
\]
फिर भी, कुछ अनुक्रम प्रणालियों में यह काम नहीं करता। अंतःप्रज्ञात्मक अनुक्रम
कलन \Log{G1i}, \Log{G1c} जैसा है, बस उसमें हर अनुक्रम की दाईं ओर अधिकतम एक
!!{formula} होने का प्रतिबंध है। यदि $\Delta = \emptyset$, तो पहला !!{proof}
\Log{G1i} में भी !!a{proof} होगा, किंतु दूसरा नहीं, क्योंकि एक अनुक्रम की
दाईं ओर $!B$ और~$\lnot !B$ दोनों हैं।

ध्यान दें कि $!A \ident \lforall[x][!B(x)]$ की स्थिति में \Log{G1c} में भी
हम नियमों का वैकल्पिक क्रम उपयोग नहीं कर सकते थे। निम्न \emph{!!a{proof}
नहीं है}:
  \[
  \AxiomC{}
  \RightLabel{$\pi_1$}
  \Deduce$!B(c) \fCenter !B(c)$
  \RightLabel{\RightR{\lforall}}
  \UnaryInf$!B(c) \fCenter \lforall[x][!B(x)]$
  \RightLabel{\LeftR{\lforall}}
  \UnaryInf$\lforall[x][!B(x)] \fCenter \lforall[x][!B(x)]$
  \DisplayProof
\]
क्योंकि~$\RightR{\lforall}$ का आइगेनचर प्रतिबंध भंग होता है।

\begin{prop}\ollabel{prop:G3c-axioms}
$!A, \Gamma \Sequent \Delta, !A$ रूप का प्रत्येक अनुक्रम ($!A$ का परमाण्विक
होना \emph{आवश्यक नहीं}) \Log{G3c} में सिद्ध किया जा सकता है।
\end{prop}

\begin{proof}
प्रमाण \olref{prop:G1c-axioms} के प्रमाण से बहुत मिलता-जुलता है, किंतु हमें
आगमन परिकल्पना में सावधान रहना होगा। $n = \depth{!A} = 0$ की स्थिति फिर सरल
है। यदि $n=0$, तो $!A$ परमाण्विक है, और $!A, \Gamma \Sequent \Delta, !A$,
\Log{G3c} की अभिगृहीति है।

अब मान लें कि $n>0$। फिर हम स्थितियाँ अलग करते हैं। आगमन परिकल्पना है: ऐसे
सभी $D!$ के लिए जिनमें $\depth{!D} < n$, और सभी $\Pi$ तथा $\Lambda$ के लिए,
$\Log{G3c} \Proves !D, \Pi \Sequent \Lambda, !D$। यहाँ $!A \ident (!B
\land !C)$ की स्थिति है। हम आगमन परिकल्पना का दो बार उपयोग करते हैं: एक बार
$!D = !B$, $\Pi = !C, \Gamma$ और $\Lambda = \Delta$ के लिए, जिससे $!B, !C,
\Gamma \Sequent \Delta, !B$ का !!a{proof}~$\pi_1$ मिलता है; और एक बार $!D
= !C$, $\Pi = !B, \Lambda$ और $\Lambda = \Delta$ के लिए, जिससे $!B, !C,
\Gamma \Sequent \Delta, !C$ का !!a{proof}~$\pi_2$ मिलता है। यह~$!B \land
!C, \Gamma \Sequent \Delta, !B \land !C$ का !!{proof} है:
\[
\AxiomC{}
\RightLabel{$\pi_1$}
\Deduce$!B, !C, \Gamma \fCenter \Delta, !B$
\RightLabel{\LeftR{\land}}
\UnaryInf$!B \land !C, \Gamma \fCenter \Delta, !B$
\AxiomC{}
\RightLabel{$\pi_2$}
\Deduce$!B, !C, \Gamma \fCenter \Delta, !C$
\RightLabel{\LeftR{\land}}
\UnaryInf$!B \land !C, \Gamma \fCenter \Delta, !C$
\RightLabel{\RightR{\land}}
\BinaryInf$!B \land !C, \Gamma \fCenter \Delta, !B \land !C$
\DisplayProof
\]

$!A \ident \lforall[x][!B(x)]$ के लिए ऐसा !!a{constant} चुनें जो
~$\lforall[x][!B(x)]$, $\Gamma$ या~$\Delta$ में न आता हो। $!D = !B(c)$,
$\Pi = \lforall[x][!B(x)], \Gamma$ और $\Lambda = \Delta$ के लिए आगमन
परिकल्पना लागू करें। इससे $!B(c), \lforall[x][!B(x)], \Gamma \Sequent
\Delta, !B(c)$ का !!a{proof}~$\pi_1$ मिलता है, जिससे हमें मिलता है:
\[
\AxiomC{}
\RightLabel{$\pi_1$}
\Deduce$!B(c), \lforall[x][!B(x)], \Gamma \fCenter \Delta, !B(c)$
\RightLabel{\LeftR{\lforall}}
\UnaryInf$\lforall[x][!B(x)], \Gamma \fCenter \Delta, !B(c)$
\RightLabel{\RightR{\lforall}}
\UnaryInf$\lforall[x][!B(x)], \Gamma \fCenter \Delta, \lforall[x][!B(x)]$
\DisplayProof
\]
जिस प्रकार हमने~$c$ चुना है, उससे $\RightR{\lforall}$ का आइगेनचर प्रतिबंध
संतुष्ट होता है।
\end{proof}

\begin{prob}
  $!A \ident (!B \lif !C)$ और $!A \ident \lexists[x][B(x)]$ की स्थितियाँ
  पूरा करके \olref{prop:G3c-axioms} का प्रमाण आगे बढ़ाएँ।
\end{prob}

\end{document}
